% =====================================================================
%  ThmSmith Stage -1 proposal skeleton.
%  Written automatically by the ThmSmith TS pipeline before Stage -1 runs.
%  Locked parameters: qid=eid\_lp\_svar\_nonequiv, specialization=v1
%
%  HOW TO USE THIS FILE
%  --------------------
%  - The preamble (everything above the `% --- BEGIN CONTENT ---` marker)
%    and the `\section{...}` headers below are fixed by the pipeline.
%    Do NOT rewrite or rename them; the Step 3 reviewer subagent and the
%    downstream Stage 0 / Stage 1 prompts grep for these section titles.
%  - Each section contains exactly one `% TODO(stage-1 ...): ...`
%    placeholder. Replace each TODO with the content described for that
%    section in `ThmSmith/tools/src/prompts/stage_neg1_draft.txt`
%    ("REQUIRED OUTPUT FORMAT (Stage -1 proposal .tex)").
%  - The `\paragraph{Locked parameters.}` block in the metadata block
%    must pin the chosen `(qid, specialization, p, assignment, target,
%    regression)` precisely. If the proposal maps to the Q1 pool-mode,
%    reproduce that block verbatim from
%    `ThmSmith/doc/panel_formula_question_generator_note_with_common_prompts.tex`
%    (Q2..Q6 templates have been retired). Otherwise (free-form
%    `(qid, specialization)` — the default), define each parameter
%    explicitly here (no implicit conventions). Stage 0 quotes this block;
%    do not rephrase it after locking.
%  - Removing a section header, renaming a macro, or rewriting the
%    preamble is a regression: the file must still match this skeleton
%    after you finish.
% =====================================================================

\documentclass[11pt]{article}

\usepackage{amsmath,amssymb,amsthm,mathtools}
\usepackage{enumitem}
\usepackage[colorlinks=true,linkcolor=blue,citecolor=blue,urlcolor=blue]{hyperref}
\usepackage{natbib}

\newcommand{\E}{\mathbb{E}}
\renewcommand{\Pr}{\mathbb{P}}
\newcommand{\one}{\mathbf{1}}
\newcommand{\transpose}{\mathsf{T}}
\newcommand{\calH}{\mathcal{H}}
\newcommand{\calF}{\mathcal{F}}
\newcommand{\calR}{\mathcal{R}}
\newcommand{\R}{\mathbb{R}}
\newcommand{\plim}{\operatorname*{plim}}

\theoremstyle{definition}
\newtheorem{definition}{Definition}
\newtheorem{assumption}{Assumption}
\newtheorem{example}{Example}

\theoremstyle{plain}
\newtheorem{theorem}{Theorem}
\newtheorem{conjecture}{Conjecture}
\newtheorem{proposition}{Proposition}
\newtheorem{lemma}{Lemma}
\newtheorem{corollary}{Corollary}

\theoremstyle{remark}
\newtheorem{remark}{Remark}

\title{ThmSmith Stage -1 proposal: \texttt{eid\_lp\_svar\_nonequiv} / \texttt{v1}%
  \\ \large\textit{Orthogonal-dependent non-Gaussian shocks as the sharp LP/SVAR non-equivalence boundary}}
\author{ThmSmith Stage -1 proposer}
\date{\today}

\begin{document}
\maketitle

% --- BEGIN CONTENT ---  (everything above is fixed by the pipeline)

\paragraph{Locked parameters.}
\begin{quote}
\begin{verbatim}
qid = eid_lp_svar_nonequiv
specialization = v1
observable distribution P = the stationary law of the reduced-form VAR
  residual u_t and lag process y_t, including Sigma_u = Var(u_t),
  the reduced-form propagation matrices Psi_h for h in H, and the fourth
  cumulant tensor K_u = cum4(u_t)
parameter theta = the identified set of structural impulse responses
  beta_hj(Q) = Psi_h Sigma_u^{1/2} Q e_j for the target shock j and
  response horizons h in H
identifying restrictions = finite-order linear VAR dynamics, orthogonal
  non-Gaussian structural shocks with finite fourth moments, Lanne-Liu
  zero co-cumulant/rank restrictions that may leave an unidentified
  orthogonal block, and a fixed partial sign/zero pattern S on the target
  shock impulse responses
bounding functional = the lower and upper support functions of the sharp
  SVAR feasible set Theta_SVAR(P,S) and of the horizonwise LP relaxation
  Theta_LP(P,S)
\end{verbatim}
\end{quote}

\begin{abstract}
Plagborg-Moller and Wolf show that local projections and VARs share the same population impulse-response object when the identifying restrictions are imposed on that common object.  Non-Gaussian SVAR work instead identifies shocks through higher moments of a single reduced-form innovation system, and recent work by Lanne, Liu, and Luoto shows that mutual independence can be relaxed to orthogonality plus skewness or kurtosis rank conditions.  This proposal asks whether the LP/SVAR equivalence survives when those rank conditions leave an unidentified orthogonal block and only partial sign/zero restrictions label the target shock.  The proposed kernel is a sharp partial-identification boundary: LP horizonwise sign feasibility is an outer relaxation of the SVAR feasible set, and strict non-equivalence appears exactly when the horizonwise sign cones in the unidentified block have no common orthogonal completion.  Resolving the conjecture would turn a verbal warning--LP rows do not encode one common shock law--into a closed-form diagnostic for applied sign-restricted non-Gaussian VAR work.
\end{abstract}

\section{Introduction}
\paragraph{Why the problem is important.}
Macroeconomists increasingly combine local projections, VAR impulse responses, sign restrictions, proxy variables, and non-Gaussian statistical identification.  The population equivalence result of \citet{PlagborgMollerWolf2021} is often read as saying that the LP-versus-VAR choice is not conceptual once the same identifying restrictions are imposed.  That reading becomes fragile when the restriction is not a reduced-form impulse-response restriction but a higher-moment restriction on one latent innovation vector: an LP coefficient at each horizon can be made sign-feasible without exhibiting one common non-Gaussian shock law across horizons.

\paragraph{The main finding (proposed).}
The headline result is Conjecture 1: under orthogonal but potentially dependent non-Gaussian shocks, the LP identified set equals the sign-restricted SVAR set if and only if the horizonwise sign/zero cones induced inside the higher-moment-unidentified orthogonal block admit a common orthogonal completion.  When the Lanne--Liu rank conditions leave an unidentified block of dimension at least two and those cones are individually feasible but jointly infeasible, the LP relaxation strictly contains structural impulse-response paths with no admissible SVAR shock law.

\paragraph{What is new.}
The novelty kernel is not another finite-sample LP/VAR comparison and not an assumption weakening of independent-component SVARs.  It is a partial-identification statement about the geometry left over after higher moments identify only a subspace: \citet{LanneLiu2026} identifies skewed or leptokurtic shocks under orthogonality, \citet{DrautzburgWright2023} refines sign-restricted sets using independence, and \citet{PlagborgMollerWolf2021} transports restrictions when the target object is common.  The epsilon-gap is the missing sharp condition for when horizonwise LP sign feasibility can be pasted into one shared orthogonal-dependent non-Gaussian structural rotation.

\paragraph{How it positions in the literature.}
The anchor cluster is partial identification for structural time-series impulse responses, with a panel/linear-projection interface.  The proposal treats the reduced-form VAR law as observed and compares two identified sets: a sharp SVAR set respecting one latent non-Gaussian shock system, and a horizonwise LP relaxation respecting the same observable propagation matrices but allowing separate rotations across horizons.  Thus the statement consumes LP/SVAR equivalence and non-Gaussian SVAR identification results as prior art, but the theorem target is a sharp-set equality/strict-containment boundary.

\paragraph{Implications for the community.}
The concrete downstream consumer is the sign-restricted independence-refinement workflow of \citet{DrautzburgWright2023}: if Conjecture 1 resolves favorably, that workflow gains a diagnostic for when LP-style horizonwise reporting is only an outer relaxation of the non-Gaussian SVAR identified set.  Researchers should report whether their sign/zero restrictions have a common completion inside the higher-moment-unidentified block before interpreting LP and SVAR sign-restricted responses as the same population object.

\section{Related work}
\citet{Jorda2005} introduced local projections as direct horizon-by-horizon impulse-response regressions, which motivates the LP side of the relaxation.  \citet{PlagborgMollerWolf2021} prove that LPs and VARs estimate the same population impulse responses under unrestricted lag structures and show that VAR identifying restrictions, including sign restrictions, can be implemented with LPs when the estimand is common.  \citet{MontielOleaPlagborgMollerQianWolf2024} show that LP inference has a double-robustness property under local lag misspecification, emphasizing that LP and VAR procedures can differ even when they are related at the population level.  \citet{LiPlagborgMollerWolf2024} document a bias--variance frontier across many DGPs and identification designs; this proposal is complementary because it targets population identified-set geometry rather than sampling performance.

\citet{GourierouxMonfortRenne2017} develop independent-component inference for SVARs and make non-Gaussianity a statistical identifying restriction on the innovation system.  \citet{LanneMeitzSaikkonen2017} show that independent non-Gaussian structural errors identify the SVAR impact matrix up to permutation and sign, producing essentially unique impulse responses.  \citet{Guay2021} gives GMM identification conditions based on higher moments and makes the rank-condition nature of non-Gaussian identification explicit.  \citet{DrautzburgWright2023} refine sign-restricted VAR sets using independence restrictions and stress weak higher-moment identification, local/non-convex identified sets, and the need to avoid assuming point identification.  \citet{LanneLiu2026} relax the independence assumption: under orthogonality and zero co-skewness or no excess co-kurtosis, skewed or leptokurtic shocks can be globally identified while the remaining shocks are set identified.

The in-catalogue prior-work summary supplied by Stage -1.1 reports no existing ThmSmith proposal about Plagborg-Moller/Wolf, SVARs, structural impulse responses, or non-Gaussian innovations.  The closest adjacent catalogue families are \texttt{q1\_spectral\_phase\_transition} on row-space frontiers, \texttt{pid\_network\_exposure} on shared-completion violations, \texttt{pid\_dose\_response\_iv} on common-latent-law pasting, and \texttt{pid\_offline\_policy\_joint\_misspec} on joint programs versus rectangular relaxations.  Those artefacts are useful algebraic analogues, but none states a time-series LP/SVAR theorem with higher-moment non-Gaussian shock restrictions.

\section{Formal setup}
Let $y_t\in\R^n$ be covariance-stationary and admit a reduced-form VAR representation with residual $u_t$ and propagation matrices $\Psi_h$, so that the reduced-form innovation contribution at horizon $h$ is $\Psi_h u_t$.  The observable law $P$ includes $\Sigma_u=\operatorname{Var}(u_t)$, $\{\Psi_h:h\in\calH\}$ for a finite horizon set $\calH$, and the fourth cumulant tensor $K_u=\operatorname{cum}_4(u_t)$.  A structural representation is $u_t=B\varepsilon_t$ with $B=\Sigma_u^{1/2}Q$, $Q\in O(n)$, $\E[\varepsilon_t]=0$, and $\operatorname{Var}(\varepsilon_t)=I_n$.  For target shock $j$, the structural response vector at horizon $h$ is
\[
  \beta_{hj}(Q)=\Psi_h \Sigma_u^{1/2}Qe_j .
\]

Let $\calS$ be a finite partial sign/zero pattern, encoded as linear inequalities and equalities $a_{\ell h}^{\transpose}\beta_{hj}(Q)\ge 0$ and $c_{\ell h}^{\transpose}\beta_{hj}(Q)=0$.  For each horizon, collect the zero-restriction coefficient rows into the matrix $D_h$.  Let $\mathcal{M}_{\mathrm{LL}}(P)$ be the set of rotations satisfying the Lanne--Liu style orthogonal-dependent higher-moment restrictions induced by $K_u$: selected co-skewness or excess co-kurtosis coordinates vanish, active marginal skewness or kurtosis coordinates are bounded away from zero, and a rank condition identifies an active shock subspace while possibly leaving an unidentified orthogonal block.  The sharp SVAR set is
\[
  \Theta_{\mathrm{SVAR}}(P,\calS)=
  \{(\beta_{hj}(Q))_{h\in\calH}: Q\in\mathcal{M}_{\mathrm{LL}}(P),\ Q\text{ satisfies }\calS\}.
\]
The LP relaxation replaces one common rotation $Q$ by horizon-specific rotations $Q_h$:
\[
  \Theta_{\mathrm{LP}}(P,\calS)=
  \{(\Psi_h\Sigma_u^{1/2}Q_he_j)_{h\in\calH}: Q_h\in\mathcal{M}_{\mathrm{LL}}(P),\ Q_h\text{ satisfies the }h\text{-part of }\calS\}.
\]
For any linear functional $\ell$ of the stacked response path, the bounding functionals are the lower and upper support values of $\ell$ over $\Theta_{\mathrm{SVAR}}(P,\calS)$ and $\Theta_{\mathrm{LP}}(P,\calS)$.

\begin{quote}
\begin{verbatim}
qid = eid_lp_svar_nonequiv
specialization = v1
observable distribution P = the stationary law of the reduced-form VAR
  residual u_t and lag process y_t, including Sigma_u = Var(u_t),
  the reduced-form propagation matrices Psi_h for h in H, and the fourth
  cumulant tensor K_u = cum4(u_t)
parameter theta = the identified set of structural impulse responses
  beta_hj(Q) = Psi_h Sigma_u^{1/2} Q e_j for the target shock j and
  response horizons h in H
identifying restrictions = finite-order linear VAR dynamics, orthogonal
  non-Gaussian structural shocks with finite fourth moments, Lanne-Liu
  zero co-cumulant/rank restrictions that may leave an unidentified
  orthogonal block, and a fixed partial sign/zero pattern S on the target
  shock impulse responses
bounding functional = the lower and upper support functions of the sharp
  SVAR feasible set Theta_SVAR(P,S) and of the horizonwise LP relaxation
  Theta_LP(P,S)
\end{verbatim}
\end{quote}

\section{Assumptions}
\begin{assumption}[Reduced-form VAR law]\label{ass:var-law}
The observable law $P$ is generated by a stable finite-order linear VAR with nonsingular residual covariance $\Sigma_u$, finite fourth moments, and known reduced-form propagation matrices $\Psi_h$ for every $h\in\calH$.
\end{assumption}

\begin{assumption}[Orthogonal-dependent non-Gaussian shock class]\label{ass:od-ng}
Admissible structural representations have $u_t=\Sigma_u^{1/2}Q\varepsilon_t$ for some $Q\in O(n)$, with $\E[\varepsilon_t]=0$, $\operatorname{Var}(\varepsilon_t)=I_n$, finite fourth moments, and no maintained mutual-independence assumption.
\end{assumption}

\begin{assumption}[Lanne--Liu rank decomposition]\label{ass:ll-rank}
The fourth-cumulant tensor $K_u$ and the imposed zero co-cumulant restrictions define an active subspace $A\subseteq\R^n$ whose shock columns are identified up to signed permutation, and an orthogonal complement $U=A^\perp$ of dimension $m\ge 0$ on which admissible rotations remain free up to $O(m)$.  The active marginal fourth-cumulant gaps are bounded below by $\delta\ge 0$; when $\delta=0$, the rank restriction is interpreted as inactive for the Gaussian-limit reduction.
\end{assumption}

\begin{assumption}[Partial sign/zero pattern]\label{ass:sign-zero}
The target shock $j$, finite horizon set $\calH$, and finite collection $\calS$ of linear sign and zero restrictions on $\{\beta_{hj}:h\in\calH\}$ are fixed before observing the higher-moment decomposition.  Restrictions may involve only the target response path and are homogeneous in the sign normalization of shock $j$.
\end{assumption}

\begin{assumption}[Nondegenerate unidentified-block exposure]\label{ass:exposure}
For every horizon $h$ used in a flagship statement, the projected loading vector $z_h=\Pi_U\Sigma_u^{1/2}\Psi_h^{\transpose}a_h$ and projected zero-restriction matrix $Z_h=\Pi_U\Sigma_u^{1/2}\Psi_h^{\transpose}D_h$ are defined by the sign/zero coefficients in Assumption~\ref{ass:sign-zero}.  At least one $z_h$ is nonzero whenever $m>0$.
\end{assumption}

\begin{assumption}[Common-completion convention]\label{ass:common-completion}
The SVAR feasible set uses one rotation $Q$ for all horizons, whereas the LP relaxation may use a separate $Q_h$ at each horizon but must use the same observable $P$, the same higher-moment admissibility rule $\mathcal{M}_{\mathrm{LL}}(P)$, and the same partial sign/zero pattern $\calS$.
\end{assumption}

\section{Main results}
\begin{theorem}[Theorem 1: common-completion inclusion and full-rank equality]\label{thm:inclusion}
Under Assumptions~\ref{ass:var-law}--\ref{ass:common-completion}, $\Theta_{\mathrm{SVAR}}(P,\calS)\subseteq\Theta_{\mathrm{LP}}(P,\calS)$.  If the Lanne--Liu rank decomposition in Assumption~\ref{ass:ll-rank} leaves no unidentified block for the target shock, i.e. $m=0$ after signed-permutation normalization of $j$, then $\Theta_{\mathrm{SVAR}}(P,\calS)=\Theta_{\mathrm{LP}}(P,\calS)$ for all homogeneous partial sign/zero patterns $\calS$.
\end{theorem}

\paragraph{Why this is new and worth studying.}
The inclusion is routine from the definitions, but the full-rank equality records the exact corner where \citet{PlagborgMollerWolf2021} and \citet{LanneMeitzSaikkonen2017} meet: the former supplies a common IRF object and the latter supplies a unique non-Gaussian shock rotation.  The useful consumer is the applied non-Gaussian SVAR workflow in \citet{LanneMeitzSaikkonen2017}: it tells users that LP reporting is harmless only after the higher-moment restrictions have removed the target-shock rotation ambiguity.

\begin{conjecture}[Conjecture 1: sharp orthogonal-block non-equivalence boundary (NOT YET PROVED)]\label{conj:orth-block-boundary}
Assume Assumptions~\ref{ass:var-law}--\ref{ass:common-completion} and suppose the Lanne--Liu decomposition leaves an unidentified block $U$ of dimension $m\ge 2$ containing the target shock.  For each horizon $h$, let
\[
  \mathcal{C}_h=\{r\in U:\|r\|=1,\ Z_h^{\transpose}r=0,\ z_h^{\transpose}r\ge 0\}
\]
be the horizon-$h$ sign/zero cone on the unidentified block.  Then the LP and SVAR sets for the target path are equal if and only if every horizonwise feasible response path admits a common $r\in\bigcap_{h\in\calH}\mathcal{C}_h$ after the signed-permutation normalization.  Equivalently, strict containment $\Theta_{\mathrm{SVAR}}(P,\calS)\subsetneq\Theta_{\mathrm{LP}}(P,\calS)$ occurs exactly on the semi-algebraic region where each $\mathcal{C}_h$ is nonempty but the common-completion intersection is empty for at least one support-function maximizer of the target response path.  In the two-dimensional unidentified-block case, the boundary is conjectured to reduce to finitely many Gram-determinant sign changes among the projected vectors $\{z_h\}$ and row spaces $\{Z_h\}$.
\end{conjecture}

\paragraph{Why this is new and worth studying.}
The closest published result is \citet{LanneLiu2026}, which characterizes which shocks are point or set identified under orthogonal non-Gaussian restrictions, but it does not compare LP horizonwise feasibility with one SVAR structural completion.  The closest LP equivalence result is \citet{PlagborgMollerWolf2021}, whose epsilon-gap is that its portability statement presumes the identifying restrictions target the same population object rather than a higher-moment latent-law completion.  Resolving this conjecture would give the \citet{DrautzburgWright2023} sign-plus-independence refinement workflow a closed-form warning: if the projected sign cones fail common completion, LP sign-restricted response paths are not structural SVAR response paths.

\begin{theorem}[Theorem 2: Gaussian and independent-component reductions]\label{thm:reductions}
Under Assumptions~\ref{ass:var-law}, \ref{ass:ll-rank}, \ref{ass:sign-zero}, and \ref{ass:common-completion}, if the fourth-cumulant gap in Assumption~\ref{ass:ll-rank} collapses to $\delta=0$ and no non-Gaussian rank restriction is imposed, the proposal reduces to the ordinary Plagborg-Moller--Wolf common-estimand comparison conditional on the same sign/zero restrictions.  If instead the shocks satisfy the independent-component restrictions of \citet{GourierouxMonfortRenne2017} or \citet{LanneMeitzSaikkonen2017} with at most one Gaussian component, the unidentified block for the target shock is empty and Theorem~\ref{thm:inclusion} gives equality.
\end{theorem}

\paragraph{Why this is new and worth studying.}
This is a sanity theorem, not the flagship contribution: it checks that the proposal does not contradict \citet{PlagborgMollerWolf2021}, \citet{GourierouxMonfortRenne2017}, or \citet{LanneMeitzSaikkonen2017} at their natural corners.  The consumer is the forthcoming double-robust LP thread of \citet{MontielOleaPlagborgMollerQianWolf2024}, because it separates population non-equivalence from the sampling and lag-misspecification issues that paper studies.

\section{Sanity-check corollaries}
\begin{corollary}[Gaussian reduced-form corner]\label{cor:gaussian}
Under Assumptions~\ref{ass:var-law}, \ref{ass:ll-rank}, \ref{ass:sign-zero}, and \ref{ass:common-completion}, if $\delta=0$ and $\mathcal{M}_{\mathrm{LL}}(P)$ imposes no higher-moment restriction, then the proposed comparison collapses to the Plagborg-Moller--Wolf common-estimand comparison with the same sign/zero restrictions, because both procedures optimize over the same second-moment rotations.
\end{corollary}

\begin{corollary}[Full non-Gaussian target identification]\label{cor:full-rank}
Under Assumptions~\ref{ass:var-law}--\ref{ass:common-completion}, if the target shock column is in the active subspace $A$ and is identified up to sign by the fourth-cumulant rank condition, then $\Theta_{\mathrm{SVAR}}(P,\calS)=\Theta_{\mathrm{LP}}(P,\calS)$ because every feasible horizon uses the same signed column after normalization.
\end{corollary}

\begin{corollary}[One-dimensional unidentified block]\label{cor:one-d}
Under Assumptions~\ref{ass:var-law}--\ref{ass:common-completion}, if $m=1$, strict LP/SVAR non-equivalence from Conjecture~\ref{conj:orth-block-boundary} cannot occur: $O(1)=\{\pm1\}$ and the sign normalization in Assumption~\ref{ass:sign-zero} leaves no continuous horizon-specific rotation to paste.
\end{corollary}

\section{Proof-sketch route}
The proof route is finite-dimensional algebra.  First whiten the reduced-form residuals and express every admissible impact matrix as $\Sigma_u^{1/2}Q$ with $Q\in O(n)$.  Next use the higher-moment rank decomposition to split $Q$ into an identified signed-permutation block and an unidentified orthogonal block $R\in O(m)$.  The SVAR set is then the image of one common $R$, whereas the LP relaxation is the product image of horizon-specific $R_h$'s.  The main work is to reduce equality of these two images to a spherical feasibility problem for the cones $\mathcal{C}_h$ and then eliminate $r$ in the $m=2$ case to obtain the claimed Gram-determinant boundary.  Boundary cases to compute explicitly are $m=0$, $m=1$, two horizons with two sign halfspaces, and a zero-restriction row that collapses a cone to antipodal points.

\section{Honest scope flags}
Theorem~\ref{thm:inclusion} and Theorem~\ref{thm:reductions} are intended as routine algebraic reductions once the feasible sets are accepted.  Conjecture~\ref{conj:orth-block-boundary} is the novel kernel and is not yet proved; the sharp ``only if'' direction and the finite Gram-determinant description in dimension two are the main open pieces.  The proposal deliberately does not claim robust inference, finite-sample coverage, or lag-misspecification validity.  The uncertain modeling choice is whether the Lanne--Liu rank decomposition should be encoded with fourth cumulants only or with a parallel skewness version; the current v1 locks the fourth-cumulant version because it matches the topic's bounded cumulant-gap language.

\section{Iteration trail}
\begin{itemize}[leftmargin=*]
\item v1 (pivot) -- initial fresh angle based on orthogonal-but-dependent non-Gaussian shocks and common-completion failure inside the Lanne--Liu unidentified block; prior verdict: none.
\end{itemize}

\section*{Bibliography}
\begin{thebibliography}{99}
\bibitem[Jord\`a(2005)]{Jorda2005}
Jord\`a, O. (2005).
Estimation and inference of impulse responses by local projections.
\emph{American Economic Review}, 95(1), 161--182.

\bibitem[Plagborg-Moller and Wolf(2021)]{PlagborgMollerWolf2021}
Plagborg-Moller, M. and Wolf, C. K. (2021).
Local projections and VARs estimate the same impulse responses.
\emph{Econometrica}, 89(2), 955--980.

\bibitem[Montiel Olea, Plagborg-Moller, Qian, and Wolf(2024)]{MontielOleaPlagborgMollerQianWolf2024}
Montiel Olea, J. L., Plagborg-Moller, M., Qian, E., and Wolf, C. K. (2024).
Double robustness of local projections and some unpleasant VARithmetic.
\emph{NBER Working Paper 32495; Econometrica forthcoming}.

\bibitem[Li, Plagborg-Moller, and Wolf(2024)]{LiPlagborgMollerWolf2024}
Li, D., Plagborg-Moller, M., and Wolf, C. K. (2024).
Local projections vs. VARs: Lessons from thousands of DGPs.
\emph{Journal of Econometrics}, 244(2), 105722.

\bibitem[Gourieroux, Monfort, and Renne(2017)]{GourierouxMonfortRenne2017}
Gourieroux, C., Monfort, A., and Renne, J.-P. (2017).
Statistical inference for independent component analysis: Application to structural VAR models.
\emph{Journal of Econometrics}, 196(1), 111--126.

\bibitem[Lanne, Meitz, and Saikkonen(2017)]{LanneMeitzSaikkonen2017}
Lanne, M., Meitz, M., and Saikkonen, P. (2017).
Identification and estimation of non-Gaussian structural vector autoregressions.
\emph{Journal of Econometrics}, 196(2), 288--304.

\bibitem[Guay(2021)]{Guay2021}
Guay, A. (2021).
GMM estimation of non-Gaussian structural vector autoregression.
\emph{Journal of Business and Economic Statistics}, 39(1), 69--81.

\bibitem[Drautzburg and Wright(2023)]{DrautzburgWright2023}
Drautzburg, T. and Wright, J. H. (2023).
Refining set-identification in VARs through independence.
\emph{Journal of Econometrics}, 235(2), 1827--1847.

\bibitem[Lanne, Liu, and Luoto(2026)]{LanneLiu2026}
Lanne, M., Liu, K., and Luoto, J. (2026).
Identifying structural vector autoregressions via non-Gaussianity of potentially dependent shocks.
\emph{The Econometrics Journal}, corrected proof.
\end{thebibliography}

\end{document}
